\scanpage{181}
\mspage{93}
% L01–04
Notons que les points de $S^{G_{s_0}}=S^{\deleted{\unsure{B}}U}$ s'identifient aux points $\dot a=aB$ \inserted{de $G/U$} dont le stabilisateur dans $G$, savoir $\operatorname{int}(a)\,\deleted{B}U$, contient $\deleted{B}U$, i.e. qui satisfont à (9).

% L04–07
\deleted{\unsure{En associant à un tel} ils correspondent donc \unsure{$\dot a$}$\in S^{G_{s_0}}$ l'unique application de $S$ dans lui-$=$}
\[
\begin{aligned}
&\deleted{f_{s_0,a}:S\longrightarrow S}\quad\deleted{\text{telle que}}\\
&\deleted{\text{1) }f_{s_0,a}\text{ commute à l'action de }G}\\
&\deleted{\text{2) }f_{s_0,a}(s_0)=a,}
\end{aligned}
\]
On trouve une bijection \deleted{\unsure{l'inverse de la bijection}} \deleted{entre l'ens.}

% L08–11
\[
\begin{aligned}
\substack{\text{ens des applications}\\\text{de }S\text{ dans lui-} =\\\text{qui commutent}\\\text{à l'action de }G}&\xrightarrow{\sim}S^{G_{s_0}}\\
f&\longmapsto f(s_0)
\end{aligned}\tag*{(11)}
\]

% L12–16
qui au niveau de l'esp. homogène quotient en \deleted{\unsure{identi|identifié}} identifiant
\[
S^{G_{s_0}}=S^{\deleted{B}U}=\{aU=\dot a\mid a\in G,\quad\operatorname{int}(a)U\supset U\tag*{(12)}
\]
donne comme bijection inverse
\[
aU=\dot a\longmapsto
\left(\substack{\text{application }\inserted{\tau_a^{d}:}\ gU\mapsto gaU\\\text{de }G/U\text{ dans lui-même}}\right)\tag*{(13)}
\]
% slanted
\margin{``translation à droite''}

% L17–23
Mais on fera attention que nous avons été obligé dans (13) de ne pas suivre la convention d'écriture (7) [où on \unsure{note|met|avait} l'exposant $-1$ en $a$] ce qui a comme conséquence qu'on aura
\[
\deleted{\unsure{f}}\tau_{ab}^{(d)}=\tau_b^{(d)}\deleted{\unsure{f}}\tau_a^{(d)}\qquad(\text{et non }\tau_a^{(d)}\tau_b^{(d)}).\tag*{(14)}
\]
Il y a de nombreux cas où l'une des

% archive sheet number: 180
